% ----------------------------------------------------------------------
\section{Optimisation continue sans contrainte}
% ----------------------------------------------------------------------

% ----------------------------------------------------------------------
\begin{frame}<beamer>
  \frametitle{Sommaire}
  \tableofcontents[currentsection]
\end{frame}

% ----------------------------------------------------------------------
\begin{frame}
  \frametitle{Enoncé du problème}

  \[
  (\text{P}^\circ) \quad
    \text{trouver} \ x \ \text{qui minimise} \ f(x), \ x \in \R^n
  \]

  \begin{itemize}
  \item problème fréquent,
  \item même avec contraintes, on peut parfois se ramener à $\text{P}^\circ$,
  \item existence d'une solution optimale si $f$ est continue et :
    \begin{itemize}
    \item soit est définie dans $S \subset \R^n$, fermé et borné,
    \item soit $f(x) \rightarrow \infty$ quand $\|x\| \rightarrow \infty$, 
    \end{itemize}
  \item mais non unicité et nombreux optima locaux pour $f$ quelconque.
  \end{itemize}
  
\end{frame}

% ----------------------------------------------------------------------
\begin{frame}
  \frametitle{En une image}

  \begin{exampleblock}{Exemple dans $\R$}
    \centering
    \includegraphics[width=0.7\textwidth,page=1]{opt-loc}
  \end{exampleblock}

\end{frame}


%% % ----------------------------------------------------------------------
%% \begin{frame}
%%   \frametitle{Optimum local}

%%   \begin{block}{Définition}
%%     La solution $x^0$ est un optimum local du problème (P) s'il existe un voisinage
%%     $V(x^0)$ autour de $x^0$ tel que $x^0$ soit optimum global du problème restreint
%%     à $V(x^0)$: 
%%   \[
%%   \left\{
%%   \begin{array}{c}
%%     \text{trouver} \ x \ \text{qui minimise} \ f(x) \ \text{où} \\
%%     \alert{x \in V(x^0)}
%%   \end{array}
%%   \right.
%%   \]
%%   \end{block}

%%   \begin{exampleblock}{Exemple dans $\R$}
%%     \centering
%%     \includegraphics[width=0.7\textwidth,page=2]{opt-loc}
%%   \end{exampleblock}
  
%% \end{frame}

% ----------------------------------------------------------------------
\begin{frame}
  \frametitle{Conditions d'optimalité générales}

  \begin{block}{Conditions d'optimalité locale}
    Soit $f$ continue et deux fois continûment différentiable en tout $x \in \R^n$.
    $x^\star$ est un optimum local ssi
    \begin{itemize}
    \item le gradient ${\nabla f}(x^\star)$ est nul (stationnarité) et
    \item le hessien $\nabla^2f(x^\star)$ est une matrice \emph{définie positive},
      c-à-d. ses valeurs propres sont \emph{strictement positives}.
      %le strictement permet d'éviter les points d'inflexion ou
      %on a un point stationnaire et ou le hessien est seulement semi-definie positive
    \end{itemize}
  \end{block}

  \only<1>{
  \begin{block}{Théorème d'optimalité globale (avec différentiabilité)}
    Soit $f$ convexe et continûment différentiable en tout $x \in \R^n$.
    $x^\star$ est un optimum global ssi ${\nabla f}(x^\star) = 0$
    (ou $f'(x^\star) = 0$ quand $n = 1$).

  \end{block}
  }
  \only<2>{
  \begin{block}{Théorème d'optimalité globale (sans différentiabilité)}
    Soit $f$ convexe et définie sur $S \subset \R^n$ convexe.
    Tout optimum local est un optimum global.  
  \end{block}
  }
  
\end{frame}

% ----------------------------------------------------------------------
\begin{frame}
  \frametitle{Exemple numérique pour $f:\R \rightarrow \R$}

  \begin{itemize}
  \item $f(x) = x^2 - 6x + 4$\\
  \item $f'(x) = 2x - 6$\\
  \item $f''(x) = 2$ \\ 
  \item $\forall x, \ f''(x) \geq 0 \Rightarrow f$ convexe
  \end{itemize}
  
  On applique le théorème : $x^\star$ est un minimum global ssi $f'(x^\star) = 0$ 
  
  On résoud l'équation : $2x^\star - 6 = 0$, donc $x^\star = 3$.
  
\end{frame}

% ----------------------------------------------------------------------
\begin{frame}
  \frametitle{Exemple numérique pour $f:\R^2 \rightarrow \R$}

  \begin{itemize}
  \item $f(x_1,x_2) = x_1^2 + x_2^2 - 4x_1 + 8x_2 - 5$
  \item ${\nabla f}(x_1,x_2) = \left(\begin{array}{l}
    2x_1 - 4\\
    2x_2 + 8
  \end{array}
    \right)$
  \item $\nabla^2f(x_1,x_2) =\left(\begin{array}{cc}
    2 & 0 \\
    0 & 2 \\
  \end{array}
    \right)$
  \item $\nabla^2f(x_1,x_2)$ définie-positive en tout $(x_1,x_2)$, 
  donc $f$ convexe
  \end{itemize}

  On applique le théorème : $x^\star$ minimum global ssi ${\nabla f}(x^\star) = 0$

  On résoud le système d'équations : 
  \[
  \left\{
  \begin{array}{ll}
    2x^\star_1 - 4 & = 0\\
    2x^\star_2 + 8 & = 0
  \end{array}
  \right.
  \]
  Donc $x^\star_1 = 2$, $x^\star_2 = -4$.
  
\end{frame}


% ----------------------------------------------------------------------
\begin{frame}
  \frametitle{Méthode numérique : descente de gradient}

  \begin{itemize}
  \item Soit $f : \R^n \rightarrow \R$ différentiable
  \item On part d'une solution initiale $x^{(0)}$
  \item A chaque étape, on obtient une nouvelle solution
    \[ x^{(k+1)} := x^{(k)} - \textcolor{blue}{\lambda^{(k)}} \alert{{\nabla f}(x^{(k)})}, \ \lambda^{(k)} > 0. \]
    (c-à-d. qu'on se déplace dans la direction opposée du \alert{gradient},
    d'un pas d'une certaine \textcolor{blue}{longueur})
  \item On s'arrête quand le gradient s'annule (ou presque)
  \end{itemize}

  ~

  Questions : 
  $x^{(0)}$ ? $\lambda^{(k)}$ ? terminaison ?
\end{frame}

% ----------------------------------------------------------------------
\begin{frame}
  \frametitle{La convergence dépend de $f$ et de $x^{(0)}$}

  \begin{itemize}
    \item $f(x) = 2x\exp{(-x^2)}$
  \end{itemize}
  
  \begin{center}
      \includegraphics[width=0.6\textwidth]{desc-grad-non-conv}    
  \end{center}
  
  \begin{itemize}
  \item En partant de $x^{(0)} = 0$, on peut atteindre $-\frac{\sqrt{2}}{2}$, le minimum.
  \item En partant de $x^{(0)} = 1$, on a $x^{(k)} \rightarrow \infty$ quand $k \rightarrow \infty$. 
  \end{itemize}
  
\end{frame}


% ----------------------------------------------------------------------
\begin{frame}
  \frametitle{Longueur du pas}

  \begin{block}{De nombreuses méthodes pour déterminer $\lambda^{(k)}$}
  \begin{itemize}
  \item fixé par l'utilisateur %(normalisé par la norme du gradient)
  \item $1 / k \|{\nabla f}(x^{(k)})\|$ (le pas diminue au fur et à mesure des itérations)
  \item \emph{Newton}: $\Big( {\nabla^2 f}(x^{(k)}) \Big)^{-1}$
    \begin{itemize}
    \item $x^{(k+1)}$ est l'optimal d'une approximation quadratique en $x^{(k)}$
    \item nécessite le hessien (et le gradient)
    \end{itemize}
  \item \emph{Steepest descent}: tel que $f(x^{(k)} - \lambda^{(k)} {\nabla f}(x^{(k)}))
    = \min_{\lambda \geq 0} f(x^{(k)} - \lambda {\nabla f}(x^{(k)}))$
  \end{itemize}
  \end{block}
\end{frame}

% ----------------------------------------------------------------------
\begin{frame}
  \frametitle{Méthode de Newton: convergence en un pas}

  \begin{itemize}
  \item $f(x_1,x_2) = a_1x_1^2 + b_1x_1 + a_2x_2^2 + b_2x_2 + c$
  \item 
    $\forall (x_1,x_2), \ {\nabla f}(x_1,x_2)^T = (2a_1x_1 + b_1,2a_2x_2 + b_2)$
  \item 
  $ \forall (x_1,x_2), \ \nabla^2f(x_1,x_2) =
  \left(\begin{array}{cc}
    2a_1 & 0 \\
    0 & 2a_2 \\
  \end{array}
  \right)
  $
  %
  \end{itemize}

  \begin{itemize}
  \item $\left( \begin{array}{c} x_1^{(1)} \\ x_2^{(1)} \end{array} \right ) =
    \left( \begin{array}{c} x_1^{(0)} \\ x_2^{(0)} \end{array} \right ) -
  \left(\begin{array}{cc} \frac{1}{2a_1} & 0 \\ 0 & \frac{1}{2a_2} \\ \end{array} \right) \cdot
  \left( \begin{array}{c} 2a_1x_1^{(0)} + b_1 \\ 2a_2x_2^{(0)} + b_2 \end{array} \right ) =
  \left( \begin{array}{c} -\frac{b_1}{2a_1} \\ -\frac{b_2}{2a_2} \end{array} \right ) $, or
  ${\nabla f}(-\frac{b_1}{2a_1}, -\frac{b_2}{2a_2})$ est nul : on a le minimum
  \item convergence en une itération car $f$ est quadratique
  \end{itemize}
  
\end{frame}

% ----------------------------------------------------------------------
\begin{frame}
  \frametitle{Méthode de Newton: non convergence}

  \begin{itemize}
  \item $f(x) = \left \{
    \begin{array}{ll}
      -4x^3 + 3x^4 & \text{si} \ x\geq 0 \\
      4x^3 + 3x^4 & \text{si} \ x < 0 \\
    \end{array} \right.$
  \end{itemize}
  
  \begin{center}
      \includegraphics[width=0.6\textwidth]{newton-non-conv}    
  \end{center}
  
  \begin{itemize}
  \item En partant de $x^{(0)} = 0.4$, on converge vers 0, le minimum.
  \item En partant de $x^{(0)} = 0.6$, on a $x^{(1)}$ = -0.6, $x^{(2)}$ = 0.6\dots 
  \end{itemize}
  
\end{frame}

% ----------------------------------------------------------------------
\begin{frame}
  \frametitle{\emph{Steepest Descent}: convergence lente}

  \begin{itemize}
  \item $f(x_1,x_2) = 2x_1x_2 + 2x_2 - x_1^2 - 2x_2^2$
  \item$(x_1^{(0)}, x_2^{(0)}) = (0, 0.5)$ 
  \end{itemize}

  \begin{center}
      \includegraphics[width=0.6\textwidth]{steepest-desc}    
  \end{center}
  
  \begin{itemize}
  \item Deux directions successives sont toujours orthogonales
  \item Ce qui est lent pour des fonctions à \emph{vallée étroite}. 
  \end{itemize}
    
  
\end{frame}

% ----------------------------------------------------------------------
\begin{frame}
  \frametitle{Autres méthodes}

  \begin{itemize}
  \item Newton modifié et quasi-Newton
  \item méthode du gradient conjugué, méthode de Fletcher et Reeves
  \item méthode de Davidon-Fletcher-Powell (DFP)
  \item méthode de Broyden-Fletcher-Goldfarb-Shanno (BFGS)
  \item algorithme de Powell (sans dérivée !)
  \item \verb!PyTorch 2.6!: Adadelta, Adafactor, Adagrad,
    Adam, AdamW, SparseAdam, Adamax, ASGD, LBFGS, NAdam,
    RAdam, RMSprop, Rprop, SGD
  \end{itemize}
  
\end{frame}

%% % ----------------------------------------------------------------------
%% \begin{frame}
%%   \frametitle{Cas d'une somme de fonctions}

%%   \[
%%   (\text{P}^\circ) \quad
%%     \text{min} \ f(x) := \sum_{i = 1}^{m} q_i(x), \ x \in \R^n
%%   \]

%%   \[ x^{(k+1)} := x^{(k)} - \lambda^{(k)} {\nabla f}(x^{(k)}), \ \lambda^{(k)} > 0, \]

%%   \[ {\nabla f}(x) = \sum_{i = 1}^{m} {\nabla q_i}(x). \]

%%   ~
  
%%   \begin{itemize}
%%   \item méthode du gradient stochastique : on évalue le gradient
%%     \begin{itemize}
%%     \item sur 1 fonction $q_i$ au hasard : ${\nabla f}(x) := {\nabla q_i}(x)$
%%     \item ou sur plusieurs (\emph{mini-batch}), mais pas sur tout $f$.
%%     \end{itemize}
%%   \end{itemize}
  
%% \end{frame}
